%! Matrices — Unit 8, Lesson 8.7 — Guided Notes (Answer Key)
\documentclass[10pt]{article}
\usepackage{precalculus-article}
\usepackage{precalculus-key}   % pulls in -boxes; do NOT also load -boxes
\newcommand{\TallMath}[1]{$\displaystyle #1\rule[-1.4em]{0pt}{3.2em}$}

\begin{document}

\pageheader{Unit 8, Lesson 8.7}{Guided Notes --- Answer Key}
\namedateperiod

\vspace{0.1in}

\begin{objectivebox}
  \textbf{Learning Targets:}
  I can identify the dimensions of a matrix and locate entries using $a_{ij}$ notation. \quad
  I can determine whether two matrices can be multiplied and compute their product.
\end{objectivebox}

\vspace{0.08in}

\begin{vocabbox}
  \textbf{matrix} --- a rectangular array of numbers arranged in rows and columns\\[3pt]
  \textbf{entry (element)} --- each individual number in a matrix; $a_{ij}$ denotes the entry in row \ans{$i$}, column \ans{$j$}\\[3pt]
  \textbf{dimensions ($n \times m$)} --- a matrix with \ans{$n$} rows and \ans{$m$} columns is called an $n \times m$ matrix\\[3pt]
  \textbf{matrix product} --- the result of multiplying two compatible matrices
\end{vocabbox}

\vspace{0.08in}

\begin{hookbox}
  \textbf{The Bakery Problem:} A bakery sells 3 items. On Day 1 it sold $[4,\,2,\,5]$ units.
  Each item costs \$3, \$5, \$2. Total revenue $=$ \ans{$4(3)+2(5)+5(2)=32$}.\\[4pt]
  \textit{Key insight: total revenue is a \ans{dot product} of the two vectors.}
\end{hookbox}

\vspace{0.08in}

\begin{notesbox}{Matrix Notation}
  \textbf{Example:} $A = \begin{bmatrix} 1 & 2 & 3 \\ 4 & 5 & 6 \end{bmatrix}$

  \begin{itemize}
    \item Dimensions of $A$: \ans{$2 \times 3$}
    \item $a_{12} = $ \ans{$2$} \quad $a_{21} = $ \ans{$4$} \quad $a_{23} = $ \ans{$6$}
  \end{itemize}

  \medskip
  \textbf{Identify dimensions and a requested entry for each matrix below:}

  \medskip
  \begin{tabular}{@{}lll@{}}
    $B = \begin{bmatrix}7 \\ 3 \\ 5\end{bmatrix}$ &
    $C = \begin{bmatrix}1 & 0 \\ 0 & 1\end{bmatrix}$ &
    $D = \begin{bmatrix}2 & 4 & 6 & 8\end{bmatrix}$\\[10pt]
    Dim: \ans{$3 \times 1$} & Dim: \ans{$2 \times 2$} & Dim: \ans{$1 \times 4$}\\[6pt]
    $b_{31} = $ \ans{$5$} & $c_{22} = $ \ans{$1$} & $d_{14} = $ \ans{$8$}
  \end{tabular}
\end{notesbox}

\vspace{0.06in}

\begin{notesbox}{Matrix Multiplication}
  \textbf{Compatibility Rule:}
  The product $AB$ is defined only when the number of \ans{columns} of $A$
  equals the number of \ans{rows} of $B$.\\[4pt]
  If $A$ is $n \times k$ and $B$ is $k \times m$, then $AB$ is \ans{$n \times m$}.

  \medskip
  \textbf{Entry Formula:}
  The $(i,j)$ entry of $AB$ is the \ans{dot product} of row \ans{$i$} of $A$
  and column \ans{$j$} of $B$.

  \medskip
  \textbf{Worked Example:}
  $\begin{bmatrix}2 & 1 \\ 0 & 3\end{bmatrix}\begin{bmatrix}4 \\ 2\end{bmatrix}$

  \medskip
  Entry $(1,1)$: row 1 of $A$ \textbf{dot} col 1 of $B$ $= 2(4) + 1(2) = $ \ans{$10$}\\[6pt]
  Entry $(2,1)$: row 2 of $A$ \textbf{dot} col 1 of $B$ $= 0(4) + 3(2) = $ \ans{$6$}

  \medskip
  Product $= $ \ans{$\begin{bmatrix}10 \\ 6\end{bmatrix}$}
\end{notesbox}

\vspace{0.06in}

\begin{practicebox}
  \textbf{Practice:} Compute each product. First check that dimensions are compatible.

  \medskip
  \begin{tabular}{@{}p{0.46\linewidth}p{0.46\linewidth}@{}}
    \textbf{1.}\quad
    $\begin{bmatrix}1 & 2 \\ 3 & 4\end{bmatrix}
     \begin{bmatrix}5 & 6 \\ 7 & 8\end{bmatrix}$

    \smallskip
    \ans{$\begin{bmatrix}1(5)+2(7) & 1(6)+2(8) \\ 3(5)+4(7) & 3(6)+4(8)\end{bmatrix} = \begin{bmatrix}19 & 22 \\ 43 & 50\end{bmatrix}$}
    &
    \textbf{2.}\quad
    $\begin{bmatrix}3 & 0 \\ 1 & 2\end{bmatrix}
     \begin{bmatrix}4 \\ 1\end{bmatrix}$

    \smallskip
    \ans{$\begin{bmatrix}3(4)+0(1) \\ 1(4)+2(1)\end{bmatrix} = \begin{bmatrix}12 \\ 6\end{bmatrix}$}
  \end{tabular}
\end{practicebox}

\end{document}
